% INSERT AFTER: online_appendix.tex lines 183-186 ("This appendix develops the probability space" to "certificates and numerical exercises.").
% INSERT BEFORE: A.1 at lines 188-189.
% NEW LABEL: oa-a-quality-class-scope. No existing label or equation number changes.
% NEW SYMBOLS: D(u)=G_L(u)-G_H(u); J_H=F_H(1), J_L=F_L(1).
% Uses the brief's j, pi, t_j, g_j, Delta_T, m(pi), M(pi), kappa(pi), and W_H.
% DEPENDENCY: update the probability-space declarations to quality classes and zero preparation cost elsewhere.
% SCOPE: A.7's displayed symmetric-signal formulas remain the equal-prior special case.
% NUMBERS: none. No new bibliography record.
% SOURCES: OA A.1-A.7; Q1_N_main Claims 2, 4, 6-8, 10-11; Q2_laplace_technical R16.

\protect\phantomsection\label{oa-a-quality-class-scope}
\noindent\textbf{Scope under quality classes.}
The class is \(j\in\{H,L\}\), with prior \(\pi=\Pr(j=H)\).
The investor observes \(j\).
Preparation reveals the value \(\theta\sim G_j\).
Conditional on \(j\), the value is independent of flow and of the preparation decision.
Finite class means imply integrable auction payoffs.
\par
The probability and deviation arguments use the class expectations
\(t_j,g_j\), the spread \(\Delta_T\), and the stated properties of Laplace noise.
They therefore apply to the quality-class model.
The point-mass formulas are special cases.
The prior and the reserve enter at the places listed below.

\begin{enumerate}
\item
In \eqref{eq:oa-oa-2} of A.1, the flow density and Bayes posterior carry the prior weights:
\[
\nu(dx)=[\pi a_H(x)+(1-\pi)a_L(x)]\,dx,\qquad
\mu_X(x)=\frac{\pi a_H(x)}{\pi a_H(x)+(1-\pi)a_L(x)}.
\]
The tower identity, the conditional laws, and the null-set argument do not change.
The likelihood ratio adds to prior log odds, so the posterior bounds become
\[
m(\pi)=\frac{\pi}{\pi+(1-\pi)e^{2/b}},\qquad
M(\pi)=\frac{\pi}{\pi+(1-\pi)e^{-2/b}}.
\]
\par
In A.3 and A.4, the residual floors are
\(\rho[1-M(\pi)]\Delta_T\) for purchases and
\(\rho m(\pi)\Delta_T\) for sales.
Their common lower bound is \(\rho\kappa(\pi)\Delta_T\), with
\(\kappa(\pi)=\min\{m(\pi),1-M(\pi)\}\).
The derivative bounds \eqref{eq:oa-oa-13} and \eqref{eq:oa-oa-15} use this constant in the strong-economy condition.
The zero preparation cost gives the entry floor at every belief.

\item
The weak-economy cost restriction in A.4, the maintained domain in A.5,
and the cost windows in A.6 use \(B_r(\pi)\).
The constant weak-economy price uses \(\pi t_H+(1-\pi)t_L\).
Under full Laplace orders,
\(\operatorname{logit}\mu_X=\operatorname{logit}\pi+(|x+1|-|x-1|)/b\), so
\[
x^*=\frac b2[\operatorname{logit}\tau-\operatorname{logit}\pi].
\]
The entry formulas after \eqref{eq:oa-oa-17}, and in A.5--A.6, use
\(\mathsf E=\pi e_H+(1-\pi)e_L\).
Ownership is \(\mathsf O_H=\pi e_HW_H\).
The upper cost restriction and the collapse condition use \(M(\pi)\).
Conditional on each class, the Laplace plateau probabilities in \eqref{eq:oa-oa-21} do not change.
Their unconditional weights are \(\pi,1-\pi\).

\item
Two further formulas in A.5 use equal priors as printed.
The no-trade existence condition becomes
\(e(\pi)\Delta_T\max\{\pi,1-\pi\}\le k\).
For the full-profile calculation \eqref{eq:oa-oa-22}, write \(J_H=F_H(1)\) and \(J_L=F_L(1)\).
Bayes' rule gives \(\pi J_H=(1-\pi)J_L\), not \(J_H=J_L\).
The sufficient existence condition is
\(k<(1-1/b)\min\{J_H,J_L\}\).
The general-prior lemma in A.6 uses the same bounds and \(\kappa(\pi)\).

\item
A.7 keeps \(\pi=1/2\) in its displayed symmetric-signal formulas.
The conditional-independence arguments apply to quality classes.
So do the conditional price equation, the measurable inverse, and the convolution argument.
At another prior, Bayes' rule gives the investor-signal probabilities and the class probabilities conditional on the signal.
Those weights enter the posterior identity \eqref{eq:oa-oa-30}, the residuals \eqref{eq:oa-oa-37}--\eqref{eq:oa-oa-38}, and the outcome probabilities after \eqref{eq:oa-oa-41}.
The buyer's private-signal maps update the public class posterior by Bayes' rule.
No claim here extends the displayed equal-prior sufficient inequalities to an arbitrary prior without these changes.

\item
The reserve enters the auction calculation in A.2 through a boundary term.
With \(D(u)=G_L(u)-G_H(u)\), the general spread is
\[
\Delta_T(F)=pF(p)D(p^-)+\int_p^\infty[1-F(u)]D(u)\,du.
\]
Equal admissibility gives \(D(p^-)=0\).
The boundary term vanishes.
Together with the rank condition above the reserve, it gives the spread comparison used in A.4--A.6.
If equal admissibility fails, the boundary term can offset the increase in the integral when the incumbent becomes stronger.
The price and deviation arguments use the resulting \(t_0,t_j,g_j,\Delta_T\).
A.7 uses \(w_j=t_j-t_0\) as well.
\end{enumerate}

The price construction in A.3 requires \(t_L-t_0\ge0\) and \(\Delta_T>0\).
Strict positivity of \(t_L-t_0\) is unnecessary.
The reserve adds no separate term to the residuals.
The common payoff level cancels.
